\documentclass[11pt]{article}
\usepackage{amsmath,amsthm,amssymb}
\usepackage{hyperref}
\usepackage[margin=1in]{geometry}

\newtheorem{theorem}{Theorem}
\newtheorem{proposition}{Proposition}
\newtheorem{corollary}{Corollary}
\newtheorem{definition}{Definition}

\title{Evidence-Bounded Inference:\\
Calibration, Abstention, and the Limits of Claim Admissibility}
\author{Roman Adriashkin\\Independent Researcher\\CogniPrint Research Initiative}
\date{Theory draft v0.1 --- November 2026}

\begin{document}
\maketitle

\begin{abstract}
We study a narrow question at the boundary of probabilistic prediction and
decision-making: when can a confidence score support a bounded claim rather
than merely a ranking? Under a strong uniform top-label calibration assumption,
we derive a selective-risk certificate for threshold abstention. We also give
an open-world impossibility result showing that no fixed non-empty acceptance
region can uniformly control false-known risk over arbitrary unknown-input
distributions, and we clarify what marginal conformal validity does and does
not imply for singleton abstention. This is a theory draft; novelty claims
remain subject to literature review and independent mathematical critique.
\end{abstract}

\section{Setup}

Let $\mathcal X$ be an observation space,
$\mathcal Y=\{1,\dots,K\}$ the known label space, and
$\mathcal Y^*=\mathcal Y\cup\{\bot\}$.
A predictor emits $p(x)\in\Delta^{K-1}$ with
\[
h(x)=\arg\max_k p_k(x),\qquad c(x)=\max_k p_k(x).
\]

\begin{definition}[Selective risk]
For an acceptance rule $g:\mathcal X\to\{0,1\}$ with
$\Pr(g(X)=1)>0$,
\[
R_{\rm sel}(g)=\Pr(h(X)\neq Y\mid g(X)=1).
\]
\end{definition}

\begin{definition}[False-known risk]
When $\Pr(Y=\bot)>0$,
\[
R_{\rm FK}(g)=\Pr(g(X)=1\mid Y=\bot).
\]
\end{definition}

Let
\[
q(x)=\Pr(Y=h(x)\mid X=x).
\]

\begin{definition}[Uniform top-label calibration defect]
The predictor has defect at most $\epsilon$ if
\[
\operatorname*{ess\,sup}_{x}|c(x)-q(x)|\le\epsilon.
\]
\end{definition}

This is substantially stronger than a binned average diagnostic such as ECE.

\section{Selective-risk certificate}

For $\tau\in[0,1]$, let $g_\tau(x)=\mathbf 1\{c(x)\ge\tau\}$.

\begin{theorem}[Threshold selective-risk certificate]
If the uniform top-label calibration defect is at most $\epsilon$ and
$\Pr(g_\tau(X)=1)>0$, then
\[
R_{\rm sel}(g_\tau)\le 1-\tau+\epsilon.
\]
\end{theorem}

\begin{proof}
On the accepted region, $c(x)\ge\tau$. Hence
$q(x)\ge c(x)-\epsilon\ge\tau-\epsilon$. Therefore the pointwise conditional
error probability is at most $1-\tau+\epsilon$. Taking the conditional
expectation over the accepted region proves the claim.
\end{proof}

\begin{corollary}[Sufficient threshold for target risk]
To certify $R_{\rm sel}\le\delta$ using the theorem, it is sufficient that
\[
\tau\ge 1-\delta+\epsilon.
\]
If the right-hand side exceeds one, this certificate is infeasible; this does
not prove that the model's actual selective risk must exceed $\delta$.
\end{corollary}

A corresponding certified correctness lower bound is
\[
s_{\rm cert}(\tau,\epsilon)=\max\{0,\tau-\epsilon\}.
\]

\begin{proposition}[Coverage bound from mean confidence]
Let $m=\mathbb E[c(X)]$ and $0\le\tau<1$. Then
\[
\Pr(c(X)\ge\tau)
\ge
\max\left\{0,\frac{m-\tau}{1-\tau}\right\}.
\]
\end{proposition}

\begin{proof}
If $C=\Pr(c(X)\ge\tau)$, then
$m\le\tau(1-C)+C=\tau+(1-\tau)C$. Rearranging gives the result.
\end{proof}


\section{Finite-sample risk certification as a baseline}

The population calibration result above assumes access to a uniform
top-label calibration defect. A finite calibration sample does not identify
that pointwise supremum without additional regularity assumptions. We
therefore separate the population theorem from an operational finite-sample
baseline.

Fix a predictor, confidence score, and threshold $\tau$ independently of an
i.i.d.\ calibration sample. Let
\[
N_\tau=\sum_{i=1}^N \mathbf 1\{c(X_i)\ge\tau\}
\]
and
\[
E_\tau=\sum_{i=1}^N
\mathbf 1\{c(X_i)\ge\tau\}\mathbf 1\{h(X_i)\neq Y_i\}.
\]
Conditional on $N_\tau=n>0$,
\[
E_\tau\sim\mathrm{Binomial}(n,R_{\rm sel}(g_\tau)).
\]
Hence a one-sided Clopper--Pearson upper confidence limit
$U_{\rm CP}(E_\tau,N_\tau;\alpha)$ gives
\[
\Pr\!\left(
R_{\rm sel}(g_\tau)
\le
U_{\rm CP}(E_\tau,N_\tau;\alpha)
\right)
\ge 1-\alpha.
\]

For a finite threshold grid
$\mathcal T=\{\tau_1,\dots,\tau_m\}$ fixed before calibration labels are
inspected, using level $\alpha/m$ for each threshold and applying the union
bound gives
\[
\Pr\!\left(
\forall j,\;
R_{\rm sel}(g_{\tau_j})\le U_j
\right)
\ge1-\alpha.
\]
A threshold may then be selected data-dependently from this fixed grid while
retaining the simultaneous high-probability statement.

This construction is a statistical baseline rather than a novelty claim.
Held-out finite-sample risk control is already developed in risk-controlling
prediction sets, Learn-then-Test, Conformal Risk Control, selective conformal
risk control, and recent selection-conditioned risk-control work. The
CogniPrint research question therefore lies above this layer: how risk
certificates should or should not authorize stronger evidence claims,
especially in the presence of unknown distributions and provenance claims.

\section{Open-world false-known impossibility}

\begin{theorem}[Arbitrary unknowns defeat any non-empty fixed acceptance region]
Let $g$ be any fixed acceptance rule with non-empty acceptance set
$A=\{x:g(x)=1\}$. If the unknown-input distribution may be arbitrary over
$\mathcal X$, then
\[
\sup_Q \Pr_{X\sim Q}(g(X)=1)=1.
\]
\end{theorem}

\begin{proof}
Choose $Q$ supported entirely on $A$.
\end{proof}

Thus a nontrivial false-known guarantee requires assumptions or evidence about
the unknown distribution; closed-set confidence alone cannot provide a
distribution-free open-world guarantee.

\section{Conformal singleton abstention}

Suppose a conformal method returns
$\Gamma_\alpha(X)\subseteq\mathcal Y$ with marginal validity
\[
\Pr(Y\in\Gamma_\alpha(X))\ge1-\alpha.
\]
Let $A=\{|\Gamma_\alpha(X)|=1\}$ and $\kappa=\Pr(A)$.

\begin{theorem}[Singleton selective-risk consequence of marginal validity]
If $\kappa>0$ and the system predicts the unique set element on $A$, then
\[
R_{\rm sel}\le \min\left\{1,\frac{\alpha}{\kappa}\right\}.
\]
\end{theorem}

\begin{proof}
A wrong accepted singleton implies
$Y\notin\Gamma_\alpha(X)$. Hence
$\Pr(\mathrm{error}\cap A)\le\alpha$. Divide by $\kappa$.
\end{proof}

Marginal conformal validity does not, by itself, guarantee
$R_{\rm sel}\le\alpha$, positive singleton coverage, or OOD rejection when the
unknown label $\bot$ is outside the conformal label space.

\section{Relation to CogniPrint}

CogniPrint M2/M2.1 motivate the theory but do not confirm it prospectively.
In particular, M2.1 showed that low false-known counts can be generated by
multi-class ambiguity rather than empty-set unknown rejection. The already
exposed Stage A records must not be used to tune a replacement method and then
present the result as confirmatory.

\section{Status and limitations}

This manuscript is a theory draft, not a peer-reviewed result. The uniform
calibration assumption is strong, finite-sample estimation of the required
defect is unresolved here, and novelty relative to selective prediction,
risk-control, calibration, and conformal inference remains to be established.

\begin{thebibliography}{9}
\bibitem{gneiting2007}
T. Gneiting and A. Raftery.
Strictly Proper Scoring Rules, Prediction, and Estimation.
\emph{JASA}, 2007.

\bibitem{guo2017}
C. Guo, G. Pleiss, Y. Sun, and K. Weinberger.
On Calibration of Modern Neural Networks.
\emph{ICML}, 2017.

\bibitem{geifman2017}
Y. Geifman and R. El-Yaniv.
Selective Classification for Deep Neural Networks.
\emph{NeurIPS}, 2017.

\bibitem{angelopoulos2023}
A. Angelopoulos and S. Bates.
Conformal Prediction: A Gentle Introduction.
\emph{Foundations and Trends in Machine Learning}, 2023.

\bibitem{clopper1934}
C. J. Clopper and E. S. Pearson.
The Use of Confidence or Fiducial Limits Illustrated in the Case of the Binomial.
\emph{Biometrika}, 1934.

\bibitem{bates2021}
S. Bates, A. Angelopoulos, L. Lei, J. Malik, and M. I. Jordan.
Distribution-Free, Risk-Controlling Prediction Sets.
arXiv:2101.02703, 2021.

\bibitem{crc2024}
A. N. Angelopoulos, S. Bates, A. Fisch, L. Lei, and T. Schuster.
Conformal Risk Control.
\emph{ICLR}, 2024.

\bibitem{lec2026}
Z. Wang et al.
LEC: Linear Expectation Constraints for Selection-Conditioned Risk Control in
Selective Prediction and Routing Systems.
\emph{ICML}, 2026.
\end{thebibliography}

\end{document}
